% PDF foreword, deliberately placed by the Pandoc template after the formal
% title page and blank verso but before the automatic table of contents.
\chapter*{Foreword}
\addcontentsline{toc}{chapter}{Foreword}

At the border, the question is rarely whether a technology can produce
an answer. The question is whether an administration can use that answer
to move a real case---faster, more consistently and with its public
authority intact. Volumes rise, fraud adapts and experienced officers
remain scarce. Those pressures make AI attractive. They also make an
unbounded AI programme dangerous.

This book is written for the leaders who must decide where to start,
which authority to retain, how to buy safely, and what evidence is
enough to scale. It does not ask an administration to trust a model or a
vendor. It asks it to begin with a bounded operational loop, measure the
result, keep the human decision where consequence requires it and expand
authority only when the evidence earns it. The central principle is
simple: decisions may transfer within a signed mandate; accountability
does not.

The evidence is deliberately uneven where the field is uneven. Mature
customs examples prove the value of data, risk analytics, document
intelligence and disciplined workforce development. They do not prove
that consequential border decisions should be handed to an autonomous
agent. That distinction is not caution for its own sake. It is the
condition that lets an ambitious administration move with credibility.

The pages that follow are therefore a book of decisions: which queue to
redesign, which rung of autonomy to grant, what must be written into the
RFP, how the system is observed in production, and when it must be
demoted. If they help one leadership team replace a vague AI ambition
with a governed, measured deployment, the book will have done its job.

\clearpage
